
\textbf{7.16.2.}
Nous ne passerons pas en revue ici, dans le cas des fibres géométriques, l'extension de toutes les propriétés des homomorphismes de spécialisation prouvées dans les sections précédentes pour les fibres «~arithmétiques~»~; elles se déduisent toutes trivialement de ces dernières par passage à la limite. Signalons seulement la conséquence intéressante suivante de 7.9.4, lorsque \(f\) est lisse et propre : \(\Gr_\bullet^{\mathrm{num}}(X_{\bar t})\) est isomorphe à un quotient d'un sous-groupe de \(\Gr_\bullet^{\mathrm{num}}(X_{\bar s})\), en particulier son rang est majoré par celui de ce dernier groupe. C'est donc un résultat de \emph{semi-continuité} pour le rang du groupe de cycles à équivalence numérique près.

\subsection*{7.17. Cas d'une base quelconque.}
Soit maintenant
\[
  f:X\longrightarrow S
\]
un morphisme plat et quasi-compact de préschémas noethériens, soient \(s,t\) deux points de \(S\) tels que \(s\) soit spécialisation de \(t\), on veut mettre en relation les groupes \(K_\bullet\) etc pour les deux fibres \(X_t\) et \(X_s\), et pour des fibres géométriques correspondantes. Pour ceci, on note (EGA II 7.1.9) qu'il existe un morphisme
\[
  g:S'\longrightarrow S,
\]
où \(S'\) est un trait, et où \(g\) transforme le point fermé \(s'\) en \(s\), le point générique \(t'\) en \(t\)~; on peut même supposer que \(k(t)\xrightarrow{\sim}k(t')\). Supposons que \(X'=X\times_S S'\) soit noethérien, ce qui est le cas si \(f\) est de type fini ou si \(S'\) est essentiellement de type fini sur \(S\) (EGA IV 1.3.8) (ce qu'on peut supposer si \(S\) lui-même est universellement japonais (EGA \(0_{\mathrm{IV}}\) 23.1.1)). Alors l'homomorphisme de spécialisation pour \(X'\) définit un homomorphisme
\begin{equation}\tag{7.17.1}
  \sigma_\bullet:K_\bullet(X_t)\xrightarrow{\sim}K_\bullet(X_{t'})
  \longrightarrow K_\bullet(X_{s'}).
\end{equation}
